\subsection{Continuous Cross-Granularity Representation Transformation}
\label{sec:coarse-support}

To model how molecular representations change across spatial granularities, we parameterize a shared continuous transformation with a neural ODE~\citep{chen2018neuralode}. We define a normalized log-granularity coordinate $u(r)=\frac{\log_2(r/8)}{4}$, which maps the training granularities $r\in\{8,32,128\}$ to $0$, $\frac{1}{2}$, and $1$, respectively, so that equal multiplicative changes in neighborhood size correspond to equal intervals.

We define the continuous transformation from the molecular representation $e_{i,a}$ at source granularity $a$ to target granularity $b$ by integrating a learned vector field $f_\theta$ over the granularity coordinate,
\begin{equation}
\frac{dz(u)}{du}=f_\theta(z(u),u),\qquad
z(u(a))=e_{i,a},\qquad T_{a\to b}(e_{i,a})=z(u(b)).
\label{eq:granularity-ode}
\end{equation}
The shared vector field defines a continuous family of cross-granularity transformations rather than separate mappings for individual granularity pairs. It is implemented as a multilayer perceptron and initialized to zero so that training starts from identity transformations. We use fourth-order Runge--Kutta integration, with implementation details provided in Appendix~\ref{app:hyperparameters}.

\subsection{Soft distributional equivariance from nested observations}
\label{sec:granularity-supervision}

Rather than enforcing exact pointwise equality between transformed and directly encoded coarse states, we constrain their relationship statistically over matched locations:
\begin{equation}
\operatorname{Law}\bigl(Z_a,T_{a\to b}(Z_a)\bigr)
\;\approx\;\operatorname{Law}(Z_a,Z_b),
\label{eq:granularity-equivariance}
\end{equation}
where $(Z_a,Z_b)$ are directly encoded states at the same sampled location. Including the source preserves fine--coarse associations: shuffling coarse states changes their pairing while preserving their marginal distribution. Two complementary losses implement this relation: coarse-support correspondence aligns predictions from fine observations within each support to a common coarse target, while JointMMD matches fine--coarse associations across locations.

\paragraph{Coarse-support correspondence.}
A coarse support supplies a common target for fine descriptions centered within it. For $B$ focal centers, define $y_{i,b}=\stopgrad(e_{i,b})$, where $\stopgrad$ stops gradients, and $\Pset=\{(8,32),(8,128),(32,128)\}$. We minimize
\begin{equation}
\mathcal L_{\mathrm{CS}}
=\frac{1}{|\Pset|B}\sum_{(a,b)\in\Pset}\sum_{i=1}^{B}
\frac{1}{bd}\sum_{j\in S_{i,b}}
\left\|T_{a\to b}(e_{j,a})-y_{i,b}\right\|_2^2.
\label{eq:coarse-support-loss}
\end{equation}
Each $e_{j,a}$ is encoded around its own center $j$ and transformed separately. This loss aligns the mean prediction to the coarse target and reduces within-support dispersion. Supports, centers and granularity pairs receive equal weight at their respective averaging levels.

\paragraph{Joint association across locations.}
For a fixed orthogonal projection $\Pi:\mathbb R^{256}\to\mathbb R^{16}$, let $\widehat P_{ab}^{T,\Pi}$ and $\widehat P_{ab}^{\Pi}$ be the empirical distributions of $(\Pi e_{i,a},\Pi T_{a\to b}(e_{i,a}))$ and $(\Pi e_{i,a},\Pi e_{i,b})$ over the same centers. Their squared maximum mean discrepancy~\citep{gretton2012mmd} gives
\begin{equation}
\mathcal L_{\mathrm{JointMMD}}=\frac{1}{|\Pset|}\sum_{(a,b)\in\Pset}
\operatorname{MMD}_{\kappa_a\otimes\kappa_b}^2
\!\left(\widehat P_{ab}^{T,\Pi},\widehat P_{ab}^{\Pi}\right).
\label{eq:granularity-joint}
\end{equation}
The product of multibandwidth radial basis kernels compares source and destination jointly. Kernel source coordinates and observed targets are detached; gradients through predicted destinations update the transition and source encoder. Appendix~\ref{app:hyperparameters} gives the estimator and correspondence-loss decomposition.

\subsection{Pretraining and continuous granularity queries}
\label{sec:training-queries}

After base pretraining, we jointly optimize the encoders and transition on the same source sections with
\begin{equation}
\mathcal L=\mathcal L_{\mathrm{base}}+\lambda_{\mathrm{CS}}\mathcal L_{\mathrm{CS}}
+\lambda_{\mathrm{J}}\mathcal L_{\mathrm{JointMMD}}.
\label{eq:full-objective}
\end{equation}
The two cross-granularity weights warm up to one. Their inputs are clean molecular encodings; detached coarse targets come from the current encoder. All supervision uses 8, 32 and 128.

After training, the encoders and transition are frozen. Retained views support tissue analysis, while $\widehat e_i(r\mid a)=T_{a\to r}(e_{i,a})$ queries trained or unseen target ranges. From a common source $a_0$, the normalized response
\begin{equation}
R_i(a,b\mid a_0)=\frac{\|\widehat e_i(b\mid a_0)-\widehat e_i(a\mid a_0)\|_2}
{\sqrt d\,\log_2(b/a)}
\label{eq:granularity-response}
\end{equation}
measures predicted change per doubling of context. We use this normalized predicted change to compare relative context sensitivity across locations.
